% ECONOMICS_PROBLEM: {"id": "D82-46", "title": "Robust dynamic mechanisms", "jel_code": "D82", "source_id": "OP-0290"}
% !TeX program = xelatex
\documentclass[11pt,a4paper]{article}
\usepackage[margin=23mm]{geometry}
\usepackage{fontspec}
\setmainfont{Latin Modern Roman}
\usepackage{amsmath,amssymb,mathtools}
\usepackage{xcolor,enumitem,booktabs,longtable,array,xurl,hyperref}
\definecolor{muted}{HTML}{53636C}
\hypersetup{colorlinks=true,urlcolor=blue,linkcolor=blue}
\setlength{\parindent}{0pt}
\setlength{\parskip}{5pt}
\newcommand{\R}{\mathbb R}
\newcommand{\E}{\mathbb E}
\newcommand{\N}{\mathbb N}
\newcommand{\ind}{\mathbf 1}
\newcommand{\Var}{\operatorname{Var}}
\newcommand{\Cov}{\operatorname{Cov}}
\newcommand{\supp}{\operatorname{supp}}
\DeclareMathOperator*{\argmin}{arg\,min}
\DeclareMathOperator*{\argmax}{arg\,max}
\newcommand{\entrymeta}[3]{{\sffamily\small\color{muted}\textbf{\texttt{#1}}\quad|\quad #2\par #3\par}}
\newcommand{\Problemref}[1]{\texttt{#1}}
\newcommand{\JELref}[2]{\texttt{#2} (JEL \texttt{#1})}
\begin{document}
\section*{Robust dynamic mechanisms}
\textbf{Problem ID:} \texttt{D82-46}\quad \textbf{JEL:} \texttt{D82}\par
% BEGIN ECONOMICS STATEMENT
\label{problem:OP-0290}
{\sffamily\small\color{muted}Original subject group: Mechanism design and market design\par}
\label{definitions:problem:30007572}
\textbf{Audit classification:} Published research agenda. Bergemann–Välimäki §8 explicitly proposes weaker informational assumptions. The displayed worst-case approximation problem is editorial; positive guarantees require an explicit domain.
\subsection*{Definitions and precise research target}
\textbf{Model and research target.} Consider a finite-horizon allocation problem with agents \(i=1,\ldots,n\), periods \(t=1,\ldots,T\), finite private type spaces \(\Theta_i\), feasible allocation set \(X\), and quasilinear discounted utility \(\sum_t\beta^{t-1}[v_i(x_t,\theta_{it})-p_{it}]\), where \(\beta\in(0,1]\). The designer knows an ambiguity set \(\mathcal P\) of joint type-transition laws, rather than one shared prior. A direct mechanism maps public report histories into allocations and transfers. Require truthful reporting to maximize every agent's expected continuation utility at every private history under every law \(P\in\mathcal P\), and require voluntary continuation under those same laws. Let \(\mathrm{Rev}_P(M)\) be expected discounted transfers and \(\mathrm{OPT}_P\) the revenue of the optimal dynamic mechanism when \(P\) is known. For a class where \(\mathrm{OPT}_P>0\), define
\[
 r^*(\mathcal P)=\sup_{M\text{ feasible and robustly truthful}}
 \inf_{P\in\mathcal P}\frac{\mathrm{Rev}_P(M)}{\mathrm{OPT}_P}.
\]
Characterize this guarantee and implementable allocation formats for ambiguity sets defined by known marginal distributions, transition bounds, or moments. Identify which restrictions permit positive guarantees and when heterogeneity of beliefs destroys the gains from intertemporal screening. Specify how zero-probability histories are handled and what information is available to each participant.

\emph{Definitions and maintained assumptions (editorial unless a source definition is named).} Specify finite type and allocation spaces, bounded valuations \(0\le v_i\le1\), payment bounds and an ambiguity set of full-support initial laws and transition kernels. Agents observe their own realized types, past allocations/transfers and declared public report history; the mechanism is a stochastic history-to-allocation/payment map with commitment. For every \(P\in\mathcal P\), every private history and every adapted reporting-and-exit strategy \(s_i\), robust dynamic IC requires \(\mathbb E_P[U_i\mid h_i,\mathrm{truth}]\ge\mathbb E_P[U_i\mid h_i,s_i]\), with opponents truthful and transition consequences of allocations included in \(P\) if applicable. Full support does not assign positive probability to every off-path report; prescribe beliefs there or impose ex-post continuation IC conditional on every possible type path. Participation must specify whether refusal is available initially or each period and how sunk payments are treated. Define \(\mathrm{OPT}_P\) using the same constraints except knowledge of \(P\). Ratio guarantees require \(\inf_P\mathrm{OPT}_P>0\), or separately handle zero benchmarks. Finite known-type optimization is an available baseline.
\subsection*{Variants and assumption changes}
\begin{enumerate}
\item \textbf{Known marginals (editorial).} Fix each agent’s one-period marginal distribution but allow all correlations and transitions in a stated finite ambiguity set. Characterize robust truthful revenue guarantees and the effect of sequential participation.
\item \textbf{Distribution-free (editorial).} Replace probabilistic continuation IC by ex-post IC for every future type path and impose bounded payments. Study approximation against the corresponding known-law optimum; this is a stronger incentive requirement with a different feasible set.
\end{enumerate}
\subsection*{References and source audit}
\begin{enumerate}
\item §8 printed pp.49–50 in June2018 revision. Supports the stated research area; exact displayed specification is editorial. DOI publication is2019; inspected revision is June2018. \url{https://cowles.yale.edu/sites/default/files/2022-09/d2102-r.pdf}
\end{enumerate}
\begin{enumerate}\raggedright
\item\hypertarget{definitions:source:30007572:1}{} Dirk Bergemann; Juuso Välimäki. 2019. \emph{Dynamic Mechanism Design: An Introduction}. Published JEL57(2),235–274 (2019); inspected June19,2018 Cowles revision §8, printed pp.49–51.
\url{https://cowles.yale.edu/sites/default/files/2022-09/d2102-r.pdf}.
DOI: \url{https://doi.org/10.1257/jel.20180892}.
\end{enumerate}

\subsection*{Common conventions: Source mathematical conventions}

These conventions apply to each entry unless explicitly replaced.

\textbf{Models and identification.} A primitive model \(m\) specifies all probability laws, feasible choices, information, equilibrium and measurement rules used by its target. The admissible class \(\mathcal M\) is fixed independently of outcomes. If \(P_O(m)\) is its observable probability law and \(\tau(m)\) the target, its identified set at observable law \(P\) is
\[
 \mathcal I_\tau(P)=\{\tau(m):m\in\mathcal M,\ P_O(m)=P\}.
\]
Point identification means that this set is a singleton. An empty set signals incompatibility of data and assumptions. Coordinate bounds need not describe the joint set. Bounds are sharp only if every retained target is attainable or arbitrarily approachable by compatible models. Finite bounds require explicit restrictions on unobserved quantities; unrestricted confounding, hidden wealth, extrapolation or arbitrary equilibrium selection can yield uninformative sets.

\textbf{Causal interventions.} A potential outcome \(Y_i(p)\) is the outcome under a complete policy regime \(p\), including timing, financing, information and market spillovers. The target population and units are fixed before treatment. With interference, use the complete assignment vector or a justified exposure mapping; individual no-interference notation alone cannot identify an economy-wide intervention. A difference of mean potential outcomes does not require the cross-world joint distribution. Conditioning on post-treatment migration, survival, enrollment or employment changes the estimand and needs separate justification.

\textbf{Statistical statements.} An estimator, confidence set, forecast or test specifies a sampling law, model class and loss/coverage criterion. Uniform \((1-\alpha)\)-coverage means \(\inf_{m\in\mathcal M}\Pr_m\{\tau(m)\in C_n\}\geq1-\alpha\), or an explicitly asymptotic version. The whole parameter space trivially has coverage, so informativeness requires a stated width, risk or power objective. A forecasting improvement is relative to a named benchmark, information set and evaluation population.

\textbf{Equilibrium and optimization.} An equilibrium comprises feasible plans, optimality, beliefs where needed, budget constraints and all market/resource clearing. The primitives or a stated benchmark define each correspondence \(\mathcal E(p;m)\). Nonemptiness is assumed or proved, never inferred just from compact policy sets. A selected-equilibrium target uses a declared selection rule; a robust target quantifies over all permitted equilibria. Use suprema or infima unless attainment is established. An \(\varepsilon\)-optimal policy is feasible and within \(\varepsilon\) of the finite optimum value. Report an empty feasible set separately.

\textbf{Welfare and accounting.} Normative utility, interpersonal normalization, social weights, numeraire, discounting and the use of public receipts are primitives. Behavioral utility need not coincide with the welfare criterion. Household welfare includes ownership income and rents once; adding profits again to compensating variation generally double-counts them. A monetary equivalent is defined by an explicit transfer and price convention, with monotonicity and range conditions. Average effects, mechanism contributions, transfers and welfare losses are different targets. Infinite sums require convergence and dynamic feasible plans require terminal or no-Ponzi conditions.

\textbf{Source versions.} A working-paper date, revision date and journal publication year are distinguished. A DOI for a journal version is not automatically the DOI of an earlier manuscript. Locators refer to the stated version and printed pages where specified. A metadata-only check does not verify a claim about a full-text conclusion. Every entry's source subsection narrows the provenance claim accordingly.
% END ECONOMICS STATEMENT
\end{document}
